\Needspace{8\baselineskip}
\subsection{Collapse of normalized Gauduchon tests}
\label{sec:gauduchon-collapse}

The root volume and the bigness threshold divide the same pairing by different quantities. In the notation $A_\theta$, $A_h$ of \cref{sec:gaud-volume-bigness-threshold}, \cref{prop:gaud-positive-volume-collapse} shows that, for $n\geq2$, a background of positive root volume that is not big forces $D_\mu(\gamma_j)\to0$ along every sequence of Gauduchon metrics with $A_h(\gamma_j)=1$ and $A_\theta(\gamma_j)\to0$, and that a positive lower bound for $D_\mu$ along one sequence minimizing the threshold would give bigness. For the radial background of \cref{thm:radial-positive-volume-nonbig}, the normalized $(n-1)$st powers of a family of balanced test metrics (\cref{def:balanced-pluriclosed}) converge weakly to the integration current of the degenerate fiber (\cref{prop:radial-localized-balanced-tests}), and the volume-minimizing family approaches the bigness threshold only in dimension three, by \cref{prop:radial-volume-test-limits}.

The quotient $D_\mu/A_h$ is always bounded above, by \eqref{eq:gaud-collapse-basic} with $\theta=h$. What is missing for the converse direction is a lower bound for $D_\mu$ along suitable sequences minimizing the threshold. Compactness of normalized potentials does not give compactness of these Gauduchon metrics.
